% =========================================================
% Proof: Uniqueness of the Derivative
% Source: volume-iii/analysis/differentiation/notes/notes-definitions.tex
% Mode: standard
% =========================================================

\subsection*{Uniqueness of the Derivative}
\label{prf:uniqueness-derivative}

\begin{remark*}[Return]
\hyperref[thm:uniqueness-derivative]{$\leftarrow$ Back to Theorem (Uniqueness of the Derivative) in Notes}
\end{remark*}

\begin{proof}
Let $I \subseteq \mathbb{R}$ be an open interval, let $f : I \to \mathbb{R}$,
and let $c \in I$. Assume $f$ is differentiable at $c$. The claim is that
$f'(c)$ is unique: if $L_1$ and $L_2$ both satisfy the definition of the
derivative at $c$, then $L_1 = L_2$.

The derivative is defined as the limit of the difference quotient at $c$.
Both $L_1$ and $L_2$ are therefore values of the same limit,
$\lim_{x \to c} \frac{f(x)-f(c)}{x-c}$. By \ByThm{Uniqueness of Limits of
Functions}, a function can have at most one limit at a cluster point.
Therefore $L_1 = L_2$.
\end{proof>

\begin{remark*}[Proof shape]
The derivative is a limit, and limits are unique. The entire proof is a
one-step reduction from uniqueness of the derivative to uniqueness of limits.

\FlashProof{Both $L_1$ and $L_2$ are values of $\lim_{x\to c}
  \frac{f(x)-f(c)}{x-c}$. By Uniqueness of Limits of Functions, this
  limit is unique, so $L_1 = L_2$.}
\end{remark*}

\begin{remark*}[Commentary]
\textbf{(a) Why is this theorem needed at all?}
The definition of the derivative says $f'(c) = L$ if $L$ satisfies
the $\varepsilon$-$\delta$ condition. Without a uniqueness theorem, there
is no guarantee that two different values $L_1$ and $L_2$ could not both
satisfy that condition. The proof closes this gap by showing the defining
limit can have at most one value — but this fact is not free, it is a
theorem about limits that must be cited.

\textbf{(b) Why does the proof reduce to limits rather than arguing directly?}
The difference quotient $\frac{f(x)-f(c)}{x-c}$ is an ordinary function
of $x$ (defined for $x \neq c$). Its limit as $x \to c$ is therefore a
function limit, and the Uniqueness of Limits of Functions theorem applies
directly. There is no need to reprove uniqueness from scratch — the limit
machinery already carries it.

\textbf{(c) What makes this proof so short?}
The work was done earlier in the curriculum when Uniqueness of Limits of
Functions was proved. This proof does nothing more than recognise that the
derivative is a special case of that theorem. Recognising when a new result
is a direct instance of a previously proved general theorem — rather than
something requiring new work — is a key proof-reading skill.
\end{remark*}

\begin{remark*}[Dependencies]
\begin{itemize}
  \item \hyperref[thm:limit-unique]{Uniqueness of Limits of Functions}:
    The single result doing all the work. The derivative is a limit;
    limits are unique; therefore the derivative is unique.
\end{itemize}

\FlashUses{thm:limit-unique}
\end{remark*}